\chapter{Conclusioni}
In questo capitolo si riassumerà lo stato di sviluppo
di \pygfa, i problemi che ancora si riscontrano e di come
sarebbe possibile risolverli. Verranno infine presentate
le conclusioni personali sull'attività di stage presentando
ciò che si è imparato nel suo sviluppo.

\section{Prima release di \pygfa}
\begin{wrapfigure} {O} {0.4\linewidth}
	\begin{centering}	
		\includegraphics[scale=0.5]{pygfa}
		\caption[Logo di \pygfa]{Logo di \pygfa.}
	\end{centering}
\end{wrapfigure}
\pygfa è alla sua prima release stabile, la copertura dei test
scritti si appresta sopra il 95\% e il numero di issue individuate
da PyLint si aggira sulle 280 ottenendo un voto dallo strumento
di 8,5/10.
Molta parte del lavoro è stata dedicata nel fornire un
ambiente che ricordasse NetworkX, sia per facilitare l'utilizzatore
sia per aiutare il futuro programmatore ad integrare nuovi moduli
della libreria nel sistema.
Il processo di sviluppo adottato ha aiutato a favorire i cambiamenti,
ad adottare strategie migliori man mano che le conoscenze aumentavano;
per questo la struttura della libreria risulta essere divisa in due metodologie
di integrazione delle funzioni di NetworkX. I primi algoritmi, riscritti
affinché venissero considerati solo archi di dovetail, sono delle
copie degli equivalenti offerti dalla libreria, modificati dei soli
nodi considerati durante le iterazioni; gli ultimi
sono stati invece pensati per essere riutilizzati da un qualsiasi iteratore
personalizzato, sfruttando la capacità del linguaggio di passare
funzioni (l'iteratore stesso) come parametri. Per questo motivo
\pygfa necessita di una fase di refactoring più profonda, che stabilisca
una struttura unica da rispettare.

L'astrazione dei dati rappresentati nel grafo è un altro problema
degno di nota, visto che in alcune situazioni genera ambiguità riguardo
le linee delle specifiche dalle quali le informazioni provengono.
Si può pensare di aggiungere un ulteriore campo agli elementi del grafo
specificando la provenienza dei dati, evitando così nelle fasi successive
di fare affidamento allo particolari configurazioni dei dati per risalire
alla linea di origine.

\section{Conoscenze acquisite}
Elencare le conoscenze che questo stage ha apportato
al mio curriculum informatico richiederebbe troppo spazio,
cercherò di riassumere le competenze acquisite che ritengo più
importanti. In primo luogo il processo di sviluppo usato: lo stage
mi ha permesso di lavorare con grande libertà per questo
un processo di sviluppo agile come extreme programming, 
che pone lo sviluppo al primo posto, mi ha senza dubbio
insegnato a gestire al meglio il tempo a disposizione e ad impostare
fin dal principio un progetto con delle fondamenta solide, ma allo stesso
tempo flessibili nell'adottare futuri cambiamenti.

Imparare e usare Python durante lo sviluppo si è rivelato più
difficile del previsto; se da una parte permette una rapida prototipazione
ed un riscontro immediato delle funzionalità, dall'altra riuscire a
capire come strutturare un progetto corposo con un linguaggio
così flessibile richiede molta pratica. Nonostante questi
dettagli, imparare ad usare questo linguaggio e gli strumenti
di test di correttezza,  della qualità del codice e della documentazione
di progetto che gli sono di corredo ha accresciuto notevolmente
non solo la mia abilità nell'implementazione di un sistema nuovo, ma anche
l'approccio col quale se ne conduce lo sviluppo.

Concluderei questo capitolo e conseguentemente questa relazione
con ciò che questo stage mi ha permesso di apprezzare più di
ogni altra cosa: lo sviluppo di sistemi performanti in grado di
lavorare con una vasta mole di informazioni quali sono le informazioni
contenute dai grafi di assemblaggio. Nonostante \pygfa non
pone molta enfasi sulle prestazioni è stato possibile
capire i contesti e le situazioni nei quali un linguaggio di alto
livello come Python risulta più indicato rispetto un linguaggio
di basso livello come il C e viceversa.





